\chapter{Commercial Relationships with Venues, Brokers and Vendors}\label{ch:fm:commercial-relationships-with-venues-brokers-and-vendors}

Nasdaq's Qualified Market Maker programme, as its November 2025 rule filing described it, paid an extra rebate of \$0.000075 a share on liquidity
added in most US stocks to a member that added more than 1.25\% of consolidated volume in a month and quoted at the national best bid or offer at
least half the time in an average of at least 2\,700 symbols a day. For a firm that already trades that much in that many names, the rebate is
money for nothing. For a firm a third of that size, meeting the terms means trading and quoting it would not otherwise do, at a cost far above the
rebate. The same published terms are a different deal for each firm, and what each can obtain in negotiation depends on what it can walk away to.

\section{What is negotiable}

A trading firm's commercial relationships are with venues (fees, rebates, programmes, connectivity), brokers (clearing, prime brokerage,
execution) and vendors (data, software, hosting). Some terms are published and the same for all: exchange fee schedules are rule filings (Book 1).
Others are negotiated: clearing rates, prime-brokerage margin and financing (chapter 14), vendor licences (chapter 22), and, within a published
framework, which programme a firm joins and how its obligations are measured.

\begin{definition}[Best alternative to a negotiated agreement, reservation value]\label{def:fm:commercial-relationships-with-venues-brokers-and-vendors:batna}
A party's \emph{best alternative to a negotiated agreement}\index{best alternative to a negotiated agreement} is what it will do if this
negotiation fails: another venue, broker or vendor, building the capability itself, or doing without. Its \emph{reservation
value}\index{reservation value} is the value of that alternative: the worst terms it should accept.
\end{definition}

\begin{definition}[Zone of possible agreement]\label{def:fm:commercial-relationships-with-venues-brokers-and-vendors:zopa}
The \emph{zone of possible agreement}\index{zone of possible agreement} is the set of terms that give each party at least its \omterm{def:fm:commercial-relationships-with-venues-brokers-and-vendors:batna}{reservation value};
when it is empty, no agreement is better for both than their alternatives.
\end{definition}

Fisher and Ury's \emph{Getting to Yes} (1981) made the best alternative the centre of practical negotiation, and Raiffa's \emph{The Art and
Science of Negotiation} (1982) its analysis. Nash (1950) gave the split that the chapter uses: with equal bargaining power, each side gets its
\omterm{def:fm:commercial-relationships-with-venues-brokers-and-vendors:batna}{reservation value} plus half of the joint gain above both \omterm{def:fm:commercial-relationships-with-venues-brokers-and-vendors:batna}{reservation values}.

\begin{proposition}[Nash's split of a transfer]\label{prop:fm:commercial-relationships-with-venues-brokers-and-vendors:nash}
Let an agreement create a gain $G$ for the venue and a cost $C$ for the firm, before a transfer $T$ from the venue to the firm, and let the
\omterm{def:fm:commercial-relationships-with-venues-brokers-and-vendors:batna}{reservation values} be $r_v$ and $r_f$. The \omterm{def:fm:commercial-relationships-with-venues-brokers-and-vendors:zopa}{zone of possible agreement} is $C+r_f\le T\le G-r_v$, non-empty exactly when $G-C\ge r_f+r_v$. The
Nash bargaining solution with the firm's bargaining weight $\beta$ is
\[
T^*=C+r_f+\beta\,(G-C-r_f-r_v).
\]
\end{proposition}

\begin{proof}
The firm's gain is $T-C$ and the venue's $G-T$; they accept when $T-C\ge r_f$ and $G-T\ge r_v$. The Nash solution maximises
$(T-C-r_f)^\beta(G-T-r_v)^{1-\beta}$; setting the derivative of its logarithm to zero gives $\beta/(T-C-r_f)=(1-\beta)/(G-T-r_v)$, whose solution is
$T^*$.
\end{proof}

\section{Market-maker programmes and tiers}

\begin{definition}[Volume commitment, shortfall penalty]\label{def:fm:commercial-relationships-with-venues-brokers-and-vendors:commit}
A \emph{volume commitment}\index{volume commitment} is a firm's undertaking to trade, add or quote at least an agreed amount on a venue or with a
counterparty over a period, in exchange for better terms. A \emph{shortfall penalty}\index{shortfall penalty} is the charge due, or the benefit
forfeited, when the commitment is not met.
\end{definition}

Book 1 introduced volume tiers and incentive programmes, and Book 11 the tier cliff: a rate that applies to every share once a threshold is met.
A market-maker programme adds obligations to the cliff, and the firm must price both: the shares it would add to reach the volume threshold (at a
loss per padded share) and the symbols it would quote at the best price that it would not quote otherwise (at a cost per symbol-day, the adverse
selection of Book 11's quoting model). \Cref{lst:fm:commercial-relationships-with-venues-brokers-and-vendors:value} values a programme this way.

\omcode{code/firm/dealterms/firm_dealterms.py}{39}{74}{A programme's daily value to a firm that pads its activity up to the requirements, the
break-even volume, the venue's gain and the zone of agreement.}\label{lst:fm:commercial-relationships-with-venues-brokers-and-vendors:value}

\begin{dated}{2026-09}{A published market-maker programme}\label{dat:fm:commercial-relationships-with-venues-brokers-and-vendors:qmm}
Nasdaq's rule filing SR-NASDAQ-2025-088 (Federal Register, 19 November 2025) describes the additional rebate of its Qualified Market Maker
programme as \$0.000075 a share in Tapes A and C (\$0.00005 in Tape B) for members adding liquidity above 1.25\% of consolidated volume, quoting at
the NBBO at least 50\% of the time in an average of 2\,700 symbols a day (1\,200 in Tape A) and having increased their added liquidity by 0.50\% of
consolidated volume since May 2020. The filing proposed replacing it with \$0.0001 a share for members adding above 0.325\% of consolidated volume
with at least 95\% of their activity adding liquidity, and a \$0.0027 credit for midpoint liquidity above 20 million shares a day. Fee programmes
change by filing.
\end{dated}

\section{Tutorial: the programme}\label{tut:fm:commercial-relationships-with-venues-brokers-and-vendors:tut}

\begin{tutorial}
\textbf{Goal.} Value a market-maker programme with published terms for firms of every size, find the volume above which it pays, and find the terms
a small and a large firm each obtain by Nash bargaining. \textbf{End state:} the value chart (\cref{fig:fm:commercial-relationships-with-venues-brokers-and-vendors:value})
and the zones of agreement (\cref{fig:fm:commercial-relationships-with-venues-brokers-and-vendors:zopa}).
\begin{tutsteps}
\item \textbf{The programme.} The dated box's terms: \$0.000075 a share on added volume, 1.25\% of consolidated volume, 2\,700 symbols at the NBBO
half the time. Consolidated volume is taken as 12 billion shares a day (an input), so the volume requirement is 150 million shares a day.
\item \textbf{The firms.} A firm adding a share $a$ of consolidated volume quotes naturally in $2\,700\,a/1\%$ symbols; padding costs \$0.001 a
share and an extra symbol at the NBBO \$2 a day (inputs). The small firm adds 0.4\% and quotes in 1\,080 symbols; the large firm adds 2\% and quotes
in 5\,400.
\item \textbf{The venue.} Each share the firm adds attracts 0.3 shares of taking flow, on which the venue earns \$0.0005 net.
\item \textbf{The bargain.} \texttt{fm\_deals.deal(firm, r\_f)} computes the zone and the Nash transfer; the large firm's outside option is a rival
venue's programme worth \$10\,000 a day, the small firm has none.
\end{tutsteps}
\end{tutorial}

\begin{omfigure}
\begin{tikzpicture}
\begin{axis}[width=10.5cm,height=5.2cm,xlabel={\small firm's natural added volume (\% of consolidated volume)},ylabel={\small net value (\$ thousand a day)},
  xmin=0,xmax=3,ymin=-140,ymax=40,tick label style={font=\footnotesize},grid=major,grid style={gray!20},table/col sep=comma]
\addplot[omDef,very thick] table[x=share,y=net]{figdata/desk/23-commercial-relationships-with-venues-brokers-and-vendors/value.csv};
\draw[black,thin] (axis cs:0,0) -- (axis cs:3,0);
\draw[gray,dotted] (axis cs:1.16,-140) -- (axis cs:1.16,40) node[pos=0.15,left,font=\scriptsize] {1.16\%};
\draw[gray,dashed] (axis cs:1.25,-140) -- (axis cs:1.25,40) node[pos=0.35,right,font=\scriptsize] {requirement 1.25\%};
\end{axis}
\end{tikzpicture}

\omcaption{The published programme's daily net value to a firm, by the firm's natural added volume: below the requirement the firm pays to pad its
volume and quote extra symbols. The value turns positive at 1.16\% of consolidated volume. Data: \texttt{fm\_deals.curve}.}\label{fig:fm:commercial-relationships-with-venues-brokers-and-vendors:value}
\end{omfigure}

The programme pays for itself only above 1.16\% of consolidated volume, just below the requirement: padding costs \$0.001 a share against a rebate
of \$0.000075, so a firm must be almost there already. For the small firm the published terms are ruinous: it earns \$11\,250 a day of rebate for
\$102\,000 of padding and \$3\,240 of extra quoting. The venue gains \$22\,500 a day from the liquidity, far less than the firm's costs: the zone
of agreement is empty, and no rebate would make the published terms worth it. A programme tailored to the small firm's natural activity (0.4\% and
1\,080 symbols) costs it nothing to meet; the venue gains \$7\,200 a day, and the Nash split gives the firm \$3\,600, which is \$0.000075 a share, the
published rate. The large firm meets the published terms without changing anything: its costs are zero, the venue gains \$36\,000 a day, and with
a rival programme worth \$10\,000 a day as its outside option the Nash split gives it \$23\,000 a day, \$0.000096 a share, 28\% above the published
rebate. Without the outside option it would get exactly the published \$18\,000 a day
(\cref{fig:fm:commercial-relationships-with-venues-brokers-and-vendors:pareto}).

\begin{omfigure}
\begin{tikzpicture}
\begin{axis}[width=10.5cm,height=4.6cm,xmin=0,xmax=110,xlabel={\small rebate paid by the venue (\$ thousand a day)},ytick={0,1,2},
  yticklabels={{small, published terms},{small, tailored terms},{large, published terms}},y dir=reverse,ymin=-0.6,ymax=2.6,
  tick label style={font=\footnotesize},grid=major,grid style={gray!20},table/col sep=comma]
\addplot[draw=none,mark=none,error bars/.cd,x dir=both,x explicit,error bar style={line width=7pt,omThm!45},error mark=none]
  table[x expr=(\thisrow{low}+\thisrow{high})/2,y=pos,x error expr=abs(\thisrow{high}-\thisrow{low})/2]{figdata/desk/23-commercial-relationships-with-venues-brokers-and-vendors/zopa.csv};
\addplot[only marks,mark=*,mark size=2.2pt,omDef] table[x=nash,y=pos]{figdata/desk/23-commercial-relationships-with-venues-brokers-and-vendors/zopa.csv};
\end{axis}
\end{tikzpicture}

\omcaption{The range between the firm's lowest acceptable rebate and the venue's highest, and the Nash rebate (dots). For the small firm under the
published terms the firm needs \$105\,000 a day and the venue gains \$22\,500: the range is inverted and there is no agreement. Data:
\texttt{fm\_deals.deal}.}\label{fig:fm:commercial-relationships-with-venues-brokers-and-vendors:zopa}
\end{omfigure}

\begin{omfigure}
\begin{tikzpicture}
\begin{axis}[width=8cm,height=6cm,xlabel={\small firm's gain (\$ thousand a day)},ylabel={\small venue's gain (\$ thousand a day)},xmin=0,xmax=40,
  ymin=0,ymax=40,tick label style={font=\footnotesize},grid=major,grid style={gray!20},table/col sep=comma,
  legend cell align=left,legend style={font=\scriptsize,at={(1.05,0.5)},anchor=west}]
\addplot[omDef,thick] table[x=firm,y=venue,restrict expr to domain={\coordindex}{0:1}]{figdata/desk/23-commercial-relationships-with-venues-brokers-and-vendors/pareto.csv};
\addplot[only marks,mark=square*,mark size=2.5pt,gray] table[x=firm,y=venue,restrict expr to domain={\coordindex}{2:2}]{figdata/desk/23-commercial-relationships-with-venues-brokers-and-vendors/pareto.csv};
\addplot[only marks,mark=*,mark size=2.8pt,omProp] table[x=firm,y=venue,restrict expr to domain={\coordindex}{3:3}]{figdata/desk/23-commercial-relationships-with-venues-brokers-and-vendors/pareto.csv};
\legend{possible splits,reservation values,Nash split}
\end{axis}
\end{tikzpicture}

\omcaption{The large firm's bargain in payoff space: every rebate splits the venue's \$36\,000 a day of gain along the line; the firm's outside
option puts its \omterm{def:fm:commercial-relationships-with-venues-brokers-and-vendors:batna}{reservation value} at \$10\,000; the Nash split halves what is left above it, \$23\,000 to the firm and \$13\,000 to the venue.
Data: \texttt{fm\_deals.deal}.}\label{fig:fm:commercial-relationships-with-venues-brokers-and-vendors:pareto}
\end{omfigure}

\section{Clearing, prime-brokerage and connectivity contracts}

The same arithmetic applies to every contract with a counterparty that values the firm's business. A clearing broker values a firm's volume and
margin balances, and its \omterm{def:fm:commercial-relationships-with-venues-brokers-and-vendors:batna}{reservation value} is the cost of the capital and operations the account uses; a prime broker values financing spreads and
balances against the \omterm{def:fm:treasury-and-funding:house}{house margin} it must hold (chapter 14); a data vendor values the licence (chapter 22). The firm's \omterm{def:fm:commercial-relationships-with-venues-brokers-and-vendors:batna}{reservation value} is its
best alternative: another broker, a second vendor, a direct connection. Each contract has terms besides price that move value: minimum
commitments with shortfall penalties, notice periods, \omterm{def:fm:the-asset-manager-and-the-fund:side}{most-favoured-nation clauses} (chapter 4), service levels (Book 14) and the right to
terminate on a change of terms.

\begin{remark}[Commitments are options written by the firm]\label{rem:fm:commercial-relationships-with-venues-brokers-and-vendors:commit}
A \omterm{def:fm:commercial-relationships-with-venues-brokers-and-vendors:commit}{volume commitment} with a \omterm{def:fm:commercial-relationships-with-venues-brokers-and-vendors:commit}{shortfall penalty} is a put the firm sells to its counterparty: if the firm's business shrinks, it pays. Its price is the
better rate; its cost is the expected penalty, largest when the firm's own volume is most uncertain, which is when a new strategy or market is
being built.
\end{remark}

\section{Leverage: what a firm of each size has}

Leverage in a negotiation is the gap between one's \omterm{def:fm:commercial-relationships-with-venues-brokers-and-vendors:batna}{reservation value} and the other side's. A large firm has outside options (rival venues compete
for its flow, several brokers want its account) and brings value the counterparty cannot easily replace; its terms improve with both. A small firm
has few alternatives and brings little that cannot be replaced; published terms designed for large firms are worth nothing to it, and its best
move is to ask for terms shaped to its activity, where the zone of agreement exists. The chapter's numbers show both: the same programme is
worth \$23\,000 a day to a large firm that bargains, and nothing to a small firm unless the terms are tailored.

\begin{method}[Negotiating a commercial contract]\label{met:fm:commercial-relationships-with-venues-brokers-and-vendors:negotiate}
\begin{enumerate}
\item Value the deal to the firm at its actual activity, including the cost of every obligation, commitment and penalty.
\item Establish the firm's best alternative and its value; improve it before negotiating (a second venue, a second broker).
\item Estimate the counterparty's gain and its alternative: what the firm's business is worth to it.
\item Negotiate terms, not only price: requirements measured on the firm's own activity, penalties capped, notice periods, termination rights.
\item Record every agreed term as data (the build of this chapter) and monitor it monthly.
\end{enumerate}
\end{method}

\section{Build: deal terms}\label{bld:fm:commercial-relationships-with-venues-brokers-and-vendors:dealterms}

\begin{build}
\bfield{Purpose} Deal packages as data, their value to both sides, \omterm{def:fm:commercial-relationships-with-venues-brokers-and-vendors:batna}{reservation values}, the zone of agreement and the Nash split.
\bfield{Interface} \texttt{firm.dealterms}: \texttt{Programme}, \texttt{Firm}, \texttt{daily\_value}, \texttt{breakeven\_share}, \texttt{venue\_value},
\texttt{zopa}, \texttt{nash\_transfer}; base fee tiers from Book 1's \texttt{firm.feesched}.
\bfield{Rules} Padding and extra quoting are priced at the firm's own costs; the rebate is a transfer; no Nash transfer when the zone is empty.
\bfield{Acceptance tests} \texttt{code/firm/dealterms/tests/}: value on hand numbers; the break-even share; the zone and the Nash split with a
bargaining weight.
\bfield{Stretch} Uncertain volume and the expected \omterm{def:fm:commercial-relationships-with-venues-brokers-and-vendors:commit}{shortfall penalty}; several venues bidding for the same flow; multi-year terms with notice.
\end{build}

\begin{omsources}
\item Nasdaq, SR-NASDAQ-2025-088, Federal Register 90 FR 52123, 19 November 2025.
\item J.~F.~Nash, ``The bargaining problem'', \emph{Econometrica} 18(2), 1950; R.~Fisher and W.~Ury, \emph{Getting to Yes}, 1981; H.~Raiffa,
\emph{The Art and Science of Negotiation}, 1982.
\end{omsources}

\section{Exercises}

\begin{exercise}[$\star$]\label{exo:fm:commercial-relationships-with-venues-brokers-and-vendors:1}
What daily rebate does a firm earn at exactly the volume requirement, and over 252 days?
\end{exercise}

\begin{exercise}[$\star$]\label{exo:fm:commercial-relationships-with-venues-brokers-and-vendors:2}
Define the \omterm{def:fm:commercial-relationships-with-venues-brokers-and-vendors:batna}{best alternative to a negotiated agreement} and the \omterm{def:fm:commercial-relationships-with-venues-brokers-and-vendors:batna}{reservation value} for a firm negotiating a clearing contract.
\end{exercise}

\begin{exercise}[$\star$]\label{exo:fm:commercial-relationships-with-venues-brokers-and-vendors:3}
What do padding and extra quoting cost the small firm each day under the published terms?
\end{exercise}

\begin{exercise}[$\star\star$]\label{exo:fm:commercial-relationships-with-venues-brokers-and-vendors:4}
Prove the Nash transfer of \cref{prop:fm:commercial-relationships-with-venues-brokers-and-vendors:nash} and apply it to the large firm.
\end{exercise}

\begin{exercise}[$\star\star$]\label{exo:fm:commercial-relationships-with-venues-brokers-and-vendors:5}
Why is the small firm's zone of agreement empty under the published terms, and what changes with tailored terms?
\end{exercise}

\begin{exercise}[$\star\star$]\label{exo:fm:commercial-relationships-with-venues-brokers-and-vendors:6}
Why is a \omterm{def:fm:commercial-relationships-with-venues-brokers-and-vendors:commit}{volume commitment} with a \omterm{def:fm:commercial-relationships-with-venues-brokers-and-vendors:commit}{shortfall penalty} dearest for a firm building a new business?
\end{exercise}

\begin{exercise}[$\star\star\star$]\label{exo:fm:commercial-relationships-with-venues-brokers-and-vendors:7}
\emph{Coding.} Give the large firm a bargaining weight of 0.3 instead of 0.5. What rebate per share does it obtain?
\end{exercise}

\begin{exercise}[$\star\star\star$]\label{exo:fm:commercial-relationships-with-venues-brokers-and-vendors:8}
\emph{Find the flaw.} ``The programme pays \$0.000075 a share: at our volume that is \$2.8 million a year, so we should join.''
\end{exercise}

\section{Problem: The Programme}

\begin{problem}[{Weekend problem --- the programme}]\label{pb:fm:commercial-relationships-with-venues-brokers-and-vendors:1}
Two firms, one small and one large, consider the same market-maker programme, and each must decide whether to join and what to ask for.

\medskip\noindent\textbf{Part I --- Negotiation.}
\begin{enumerate}
\item Define the best alternative, the \omterm{def:fm:commercial-relationships-with-venues-brokers-and-vendors:batna}{reservation value} and the \omterm{def:fm:commercial-relationships-with-venues-brokers-and-vendors:zopa}{zone of possible agreement}.
\item State and prove \cref{prop:fm:commercial-relationships-with-venues-brokers-and-vendors:nash}.
\item What is negotiable with a venue, a broker and a vendor?
\item Define a \omterm{def:fm:commercial-relationships-with-venues-brokers-and-vendors:commit}{volume commitment} and a \omterm{def:fm:commercial-relationships-with-venues-brokers-and-vendors:commit}{shortfall penalty}.
\end{enumerate}

\medskip\noindent\textbf{Part II --- The programme.}
\begin{enumerate}[resume]
\item State the published terms and their source.
\item How does the chapter cost the obligations?
\item Give the net value curve and the break-even volume.
\item Why is the break-even so close to the requirement?
\end{enumerate}

\medskip\noindent\textbf{Part III --- The bargains.}
\begin{enumerate}[resume]
\item Give the small firm's value, the venue's gain and the zone under the published terms.
\item Give the tailored terms' zone and Nash transfer.
\item Give the large firm's zone and Nash transfer, with and without its outside option.
\item What does the outside option add?
\item How would a lower bargaining weight change the large firm's rebate?
\end{enumerate}

\medskip\noindent\textbf{Part IV --- Decisions.}
\begin{enumerate}[resume]
\item Should the small firm join? What should it ask for?
\item Should the large firm join? What should it ask for?
\item How would you improve the small firm's best alternative?
\item Which non-price terms would you negotiate?
\item How would you monitor the agreement once signed?
\item State the \emph{named result}: the volume above which a market-maker programme's discount pays for its obligations, and the Nash-bargaining
fee that a small and a large firm each obtain.
\item In two sentences, write each firm's negotiating position.
\end{enumerate}
\end{problem}

\section{Interview questions}

\begin{interviewq}[$\star$ \iqroles{trader}]\label{iq:fm:commercial-relationships-with-venues-brokers-and-vendors:1}
What is a market-maker programme, and what does a venue get from it?
\end{interviewq}

\begin{interviewq}[$\star$ \iqroles{trader}]\label{iq:fm:commercial-relationships-with-venues-brokers-and-vendors:2}
What is your best alternative when negotiating fees with your clearing broker?
\end{interviewq}

\begin{interviewq}[$\star\star$ \iqroles{researcher}]\label{iq:fm:commercial-relationships-with-venues-brokers-and-vendors:3}
A rebate programme requires a volume you reach in only half the months. How do you value it?
\end{interviewq}

\begin{interviewq}[$\star\star$ \iqroles{trader, researcher}]\label{iq:fm:commercial-relationships-with-venues-brokers-and-vendors:4}
Derive the Nash bargaining split of a surplus between two parties with outside options.
\end{interviewq}

\begin{interviewq}[$\star\star$ \iqroles{risk}]\label{iq:fm:commercial-relationships-with-venues-brokers-and-vendors:5}
What risks does a \omterm{def:fm:commercial-relationships-with-venues-brokers-and-vendors:commit}{volume commitment} create?
\end{interviewq}

\begin{interviewq}[$\star\star\star$ \iqroles{researcher, trader}]\label{iq:fm:commercial-relationships-with-venues-brokers-and-vendors:6}
Two venues compete for your flow with rebate programmes. How would you decide where to send it?
\end{interviewq}
